% !TeX program = pdflatex
%==============================================================================
%  Lista Teórica 10 --- KNN e Árvores de Classificação
%  O gabarito sai de "Lista teorica 10 - gabarito.tex".
%==============================================================================
\providecommand{\gabaritoopt}{}
\documentclass[11pt,a4paper]{article}
\usepackage[\gabaritoopt]{../../recursos/latex/estilo-lista}

\begin{document}
\cabecalholista{10}{KNN e Árvores de Classificação}

%------------------------------------------------------------------------------
\begin{exercicio}
Uma engenheira de qualidade de uma fábrica de cerâmica usa uma árvore de
classificação para prever se um lote de peças será reprovado ($Y=1$) ou aprovado
($Y=0$) na inspeção final, a partir de medidas do processo de produção. Para uma região $R$ do espaço das covariáveis, sejam $n_R$ o número
de lotes de treino em $R$, $n_{R,c}$ o número deles na classe $c\in\{0,1\}$ e
$\widehat p_{R,c}=n_{R,c}/n_R$. O erro de classificação, o índice de Gini e a
entropia de $R$ são
\[
  E(R)=1-\max_{c}\,\widehat p_{R,c},\qquad
  G(R)=\sum_{c}\widehat p_{R,c}\,(1-\widehat p_{R,c}),\qquad
  D(R)=-\sum_{c}\widehat p_{R,c}\log_2\widehat p_{R,c},
\]
com a convenção $0\log_20=0$. Se $I$ é uma das três medidas e $R$ é dividida em $R_1$
e $R_2$, com $n_1$ e $n_2$ lotes, a impureza ponderada do corte é
$n_1I(R_1)+n_2I(R_2)$, e a redução de impureza é
$n_RI(R)-\big[n_1I(R_1)+n_2I(R_2)\big]$. Como $n_RI(R)$ não depende do corte, o corte
de menor impureza ponderada é o de maior redução.

Durante o crescimento da árvore, a engenheira examina uma região $R$ com $180$ lotes e
um corte pela umidade da argila, que divide $R$ em $R_1$ e $R_2$:
\begin{center}\small
\renewcommand{\arraystretch}{1.25}
\begin{tabular}{@{}lccc@{}}
\toprule
região & lotes & reprovados & aprovados \\
\midrule
$R$   & $180$ & $30$ & $150$ \\
$R_1$ & $90$  & $0$  & $90$ \\
$R_2$ & $90$  & $30$ & $60$ \\
\bottomrule
\end{tabular}
\end{center}
\begin{enumerate}[label=(\alph*),leftmargin=1.1cm]
  \item Com duas classes, escreva $p=\widehat p_{R,1}$ e mostre que
        $E(R)=\min\{p,1-p\}$, $G(R)=2p(1-p)$ e
        $D(R)=-p\log_2p-(1-p)\log_2(1-p)$. Mostre também que $n_RE(R)$ é o número de
        lotes de treino de $R$ que a moda de $R$ classifica errado.
  \item Calcule, para cada uma das três medidas, $n_RI(R)$, a impureza ponderada
        $n_1I(R_1)+n_2I(R_2)$ e a redução de impureza do corte da engenheira. Use
        $\log_23\approx1{,}585$ e $\log_25\approx2{,}322$.
  \item O corte é bom? Qual das três medidas não registra isso? Explique o valor que
        ela dá contando os erros de treino da moda em $R$, em $R_1$ e em $R_2$.
  \item Considere também um segundo corte, por outra covariável, que divide $R$ em duas
        regiões de $90$ lotes, cada uma com $15$ reprovados. Calcule as três reduções
        desse corte e compare os dois cortes por cada uma das medidas. O que
        aconteceria se o algoritmo que cresce a árvore usasse o erro de classificação?
\end{enumerate}
\end{exercicio}

\begin{solucao}
\textbf{(a)} Com duas classes, $\widehat p_{R,1}=p$ e $\widehat p_{R,0}=1-p$, porque
cada lote de $R$ pertence a exatamente uma delas. Como $1-\max\{a,b\}=\min\{1-a,1-b\}$,
\[
  E(R)=1-\max\{p,1-p\}=\min\{1-p,\,p\},
\]
e, somando as parcelas das duas classes,
\[
  G(R)=p\,(1-p)+(1-p)\,\big(1-(1-p)\big)=2p(1-p).
\]
A expressão de $D(R)$ é a própria definição, escrita com as suas duas parcelas.

Se $R$ for uma folha, a árvore prediz para todos os lotes de $R$ a sua moda, uma classe
$c^*$ com $n_{R,c^*}=\max_cn_{R,c}$. Ela acerta os $n_{R,c^*}$ lotes dessa classe e erra
todos os outros, que são
\[
  n_R-\max_cn_{R,c}=n_R\Big(1-\max_c\widehat p_{R,c}\Big)=n_RE(R).
\]
Se houver empate entre as classes, o número de erros é o mesmo qualquer que seja a
classe escolhida. Assim, $E(R)$ é a taxa de erro de treino da moda de $R$, e $n_RE(R)$
é o número desses erros.

\smallskip
\textbf{(b)} A proporção de reprovados é $30/180=1/6$ em $R$, $0$ em $R_1$ e
$30/90=1/3$ em $R_2$. A região $R_1$ é pura, e as suas três impurezas são nulas: o
item (a) dá $E(R_1)=G(R_1)=0$ e, com a convenção $0\log_20=0$, $D(R_1)=0$. Nas outras
duas regiões,
\[
  G(R)=2\cdot\tfrac16\cdot\tfrac56=\tfrac5{18},
  \qquad
  G(R_2)=2\cdot\tfrac13\cdot\tfrac23=\tfrac49,
\]
e, como $-\log_2p=\log_2(1/p)$,
\[
  D(R)=\tfrac16\log_26+\tfrac56\log_2\tfrac65=\log_26-\tfrac56\log_25
       =1+\log_23-\tfrac56\log_25\approx0{,}650,
\]
\[
  D(R_2)=\tfrac13\log_23+\tfrac23\log_2\tfrac32=\log_23-\tfrac23\approx0{,}918 .
\]
Multiplicando cada impureza pelo número de lotes da região:
\begin{center}\small
\renewcommand{\arraystretch}{1.4}
\begin{tabular}{@{}lccc@{}}
\toprule
medida & $n_RI(R)$ & $n_1I(R_1)+n_2I(R_2)$ & redução \\
\midrule
$E$ & $180\cdot\frac16=30$ & $90\cdot0+90\cdot\frac13=30$ & $0$ \\
$G$ & $180\cdot\frac5{18}=50$ & $90\cdot0+90\cdot\frac49=40$ & $10$ \\
$D$ & $180\cdot0{,}650\approx117{,}0$ & $90\cdot0+90\cdot0{,}918\approx82{,}6$ & $\approx34{,}4$ \\
\bottomrule
\end{tabular}
\end{center}
Em forma exata, a redução da entropia é $240+90\log_23-150\log_25\approx34{,}36$.

As reduções de medidas diferentes não se comparam entre si, porque cada medida tem a
sua escala: com duas classes, $E$ e $G$ valem no máximo $1/2$, e $D$, no máximo $1$.
O que se compara são cortes diferentes pela mesma medida, como no item (d).

\smallskip
\textbf{(c)} O corte é bom. A região $R_1$ é pura: os $90$ lotes que ela contém, metade
dos lotes de $R$, foram todos aprovados. Para um lote novo com umidade abaixo do corte,
os dados de treino não registram nenhuma reprovação; antes do corte, a mesma pergunta
tinha como resposta a proporção de $R$, um lote reprovado em cada seis. As $30$
reprovações ficam concentradas em $R_2$.

A medida que não registra isso é o erro de classificação, cuja redução é exatamente
zero. Pelo item (a), $n_RE(R)$, $n_1E(R_1)$ e $n_2E(R_2)$ são os números de erros de
treino das modas das três regiões. A moda é ``aprovado'' em todas elas: $150$ contra
$30$ em $R$, $90$ contra $0$ em $R_1$ e $60$ contra $30$ em $R_2$. Antes do corte, a
moda de $R$ erra os $30$ lotes reprovados; depois, a de $R_1$ não erra nenhum lote, e a
de $R_2$ erra os mesmos $30$. O corte não muda a classe prevista para nenhum lote de
treino, e por isso não pode mudar o número de erros, que é tudo o que $E$ mede.

O Gini e a entropia, ao contrário, registram o ganho. Para ver de onde ele vem, use
$n_R=n_1+n_2$ e escreva a redução região por região:
\[
  n_RI(R)-\big[n_1I(R_1)+n_2I(R_2)\big]
  =n_1\big[I(R)-I(R_1)\big]-n_2\big[I(R_2)-I(R)\big].
\]
O primeiro termo é o ganho de pureza dos lotes que foram para $R_1$, e o segundo, a
perda dos que foram para $R_2$, onde a proporção de reprovados subiu de $1/6$ para
$1/3$. As três medidas concordam que $R_1$ melhorou e $R_2$ piorou; o que muda é o peso
de cada parte. Para o erro de classificação,
$90\big(\tfrac16-0\big)-90\big(\tfrac13-\tfrac16\big)=15-15=0$, e o ganho e a perda se
cancelam. Para o Gini, $90\big(\tfrac5{18}-0\big)-90\big(\tfrac49-\tfrac5{18}\big)=25-15=10$,
e, para a entropia, $90\,(0{,}650-0)-90\,(0{,}918-0{,}650)\approx58{,}5-24{,}1=34{,}4$:
o ganho supera a perda.

É a mesma situação do exemplo dos dados \texttt{Heart} no [ISLP] §8.1.2: um corte cujas
duas folhas preveem a mesma classe não reduz o erro de classificação e, ainda assim, é
feito, porque aumenta a pureza.

\smallskip
\textbf{(d)} No segundo corte, as duas regiões têm a mesma proporção de reprovados que
$R$, $15/90=1/6$. Como as três medidas dependem da região só pelas proporções das
classes, cada região tem a mesma impureza que $R$, e
\[
  n_1I(R_1)+n_2I(R_2)=90\,I(R)+90\,I(R)=n_RI(R):
\]
as três reduções são nulas. O corte não separa nada, e as duas regiões reproduzem, em
proporção, a região original.

Pelo Gini e pela entropia, os dois cortes ficam bem distintos, com reduções $10$ contra
$0$ e $34{,}4$ contra $0$. Pelo erro de classificação, os dois têm redução zero e ficam
empatados. Se o algoritmo que cresce a árvore usasse o erro de classificação, ele não
teria como preferir o corte da engenheira, que isola $90$ lotes sem nenhuma reprovação,
a um corte que não muda nada. O argumento do item (c) mostra mais: sempre que as duas
partes de um corte mantêm a classe prevista em $R$, nenhuma previsão de treino muda, e a
redução do erro de classificação é zero. Numa região em que todos os cortes candidatos
fossem desse tipo, esse critério não distinguiria corte nenhum. É por isso que a árvore
cresce com o Gini ou com a entropia.
\end{solucao}

%------------------------------------------------------------------------------
\begin{exercicio}
Um estudante quer entender o papel de $k$ no KNN de classificação e começa pelos dois
valores extremos, $k=1$ e $k=n$. Ele tem $n=301$ observações de treino de duas classes,
$181$ da classe~$1$ e $120$ da classe~$0$, com vetores de covariáveis
$\X_1,\dots,\X_n$ distintos dois a dois, e usa a distância euclidiana $d$. Lembre-se que o
classificador é
\[
  \ghat(\x)=\operatorname{moda}\{Y_i:i\in\mathcal{N}_{\x}^{(k)}\},
  \qquad\text{com}\qquad
  \mathcal{N}_{\x}^{(k)}=\{i:d(\X_i,\x)\le d^k_{\x}\},
\]
em que $d^k_{\x}$ é a distância de $\x$ à $k$-ésima observação de treino mais próxima.
\begin{enumerate}[label=(\alph*),leftmargin=1.1cm]
  \item Mostre que, com $k=1$, o classificador acerta todas as observações de treino.
        Mostre que, com $k=n$, ele atribui a mesma classe a todo $\x$. Calcule o erro de treino
        para cada um desses casos.
  \item Descreva a fronteira de decisão nos dois casos. Suponha, em seguida, que o
        estudante descubra que alguns rótulos foram registrados errado e os corrija,
        trocando $1$ por $0$. Qual é o menor número de rótulos trocados que muda o
        classificador com $k=1$? E com $k=n$? Relacione as respostas com o viés e a
        variância de cada caso.
  \item Pelo item (a), o estudante conclui que $k=1$ é a melhor escolha, porque com ela
        o classificador não erra nenhuma observação de treino. O argumento está certo?
        Como $k$ deve ser escolhido, e por que esse procedimento não teria o mesmo problema de $k=1$?
\end{enumerate}
\end{exercicio}

\begin{solucao}
\textbf{(a)} Com $k=1$, tome uma observação de treino $\X_i$. Como $d$ é uma distância,
$d(\X_i,\X_i)=0$, e $d(\X_j,\X_i)>0$ para todo $j\ne i$, porque os vetores são
distintos. Logo $d^1_{\X_i}=0$, $\mathcal{N}^{(1)}_{\X_i}=\{i\}$ e
$\ghat(\X_i)=\operatorname{moda}\{Y_i\}=Y_i$: cada observação de treino é a sua própria
vizinha mais próxima, e a taxa de erro de treino é zero. A hipótese de vetores distintos
é necessária. Se $\X_i=\X_j$ com $Y_i\ne Y_j$, o classificador atribui uma única classe a
esse ponto e erra ao menos uma das duas observações, qualquer que seja a regra de
desempate.

Com $k=n$, $d^n_{\x}$ é a distância de $\x$ à observação de treino mais distante, e
todas as observações satisfazem $d(\X_i,\x)\le d^n_{\x}$. Então
$\mathcal{N}^{(n)}_{\x}=\{1,\dots,n\}$ para todo $\x$, e $\ghat(\x)$ é a moda de todos
os rótulos de treino: a classe~$1$, com $181$ observações contra $120$. O classificador
não depende de $\x$; é o classificador constante que prediz a classe majoritária. Ele
erra exatamente as $120$ observações da classe~$0$, e a taxa de erro de treino é
$120/301\approx0{,}399$, a proporção da classe minoritária.

\smallskip
\textbf{(b)} Com $k=1$, cada $\x$ recebe a classe da observação de treino mais próxima.
O espaço fica dividido em regiões, uma para cada $\X_i$, formadas pelos pontos mais
próximos de $\X_i$ do que de qualquer outra observação de treino, e a fronteira de
decisão é formada pelas bordas comuns a regiões de observações de classes diferentes.
Ela acompanha cada observação: é muito recortada, e uma observação isolada no meio da
outra classe cria uma ilha em torno de si, como na Figura~1 da Aula~10. Com $k=n$, não
há fronteira: todo $\x$ é atribuído à classe~$1$.

Com $k=1$, um único rótulo trocado já muda o classificador. Se $Y_i$ passa de $1$ a $0$,
o item (a) dá $\ghat(\X_i)=0$, onde antes era $1$. A mudança não se restringe ao ponto
$\X_i$. Seja $\delta=\min_{j\ne i}d(\X_i,\X_j)$, que é positivo porque os vetores são
distintos. Se $d(\x,\X_i)<\delta/2$, a desigualdade triangular dá, para todo $j\ne i$,
\[
  d(\x,\X_j)\ \ge\ d(\X_i,\X_j)-d(\x,\X_i)\ >\ \delta-\frac{\delta}{2}
  \ =\ \frac{\delta}{2}\ >\ d(\x,\X_i),
\]
de modo que $\X_i$ é a vizinha mais próxima de $\x$, e a previsão muda em toda a bola de
raio $\delta/2$ centrada em $\X_i$.

Com $k=n$, o classificador só muda se mudar a classe majoritária. Com $m$ rótulos
trocados, as classes passam a ter $181-m$ e $120+m$ observações, e a classe~$0$ se torna
majoritária se e somente se $120+m>181-m$, isto é, $m>30{,}5$. Com até $30$ trocas, o
classificador não muda em ponto nenhum; com $31$, quaisquer que sejam, muda em todos.

Em classificação, a aula usa viés e variância no mesmo sentido da regressão: a variância
diz o quanto o classificador ajustado muda quando a amostra de treino muda, e o viés, o
quanto ele erra de forma sistemática, mesmo com muitos dados. As contagens acima medem a
variância de forma direta. Com $k=1$, o classificador depende de cada observação
individualmente, e um rótulo diferente basta para mudá-lo numa região inteira: a
variância é alta. Com $k=n$, ele depende da amostra só pela classe majoritária, que
resiste a $30$ trocas em $301$: a variância é baixa. O viés vai no sentido oposto. Com
$k=n$, o classificador ignora $\x$; se a proporção de cada classe varia ao longo do
espaço das covariáveis, ele erra sistematicamente onde a classe~$0$ predomina, por maior
que seja $n$, e o viés é alto. Com $k=1$, a decisão em $\x$ vem da observação de treino
mais próxima, e, quando há observações perto de $\x$, ela acompanha essa variação: o viés
é baixo.

\smallskip
\textbf{(c)} O argumento não está certo. O item (a) mostrou que o erro de treino do
$1$-NN é zero em qualquer conjunto de treino com vetores distintos, qualquer que seja a
relação entre $\X$ e $Y$. Ele seria zero mesmo que $Y$ fosse independente de $\X$. Nesse
caso, as covariáveis não trazem informação nenhuma sobre a classe, o classificador de
Bayes da Aula~07 prediz sempre a classe mais provável, e nenhum classificador erra menos
do que ele em observações novas. Um número que vale zero em qualquer situação não informa
nada sobre o desempenho em observações novas, que é o que o risco mede. O erro de treino
do $1$-NN é zero porque cada observação de treino é usada para prever a si mesma.

Como a aula indica, $k$ deve ser escolhido por validação cruzada. Nela, cada observação
é prevista por um classificador ajustado sem ela. No LOOCV com $k=1$, por exemplo, a
observação $i$ é prevista pelo rótulo da mais próxima entre as outras $n-1$, e há erro
sempre que esse rótulo é diferente de $Y_i$. Como a observação deixada de fora não é
mais vizinha de si mesma, todos os valores de $k$ passam a ser comparados pelo erro em
observações que não participaram do ajuste, que estima o risco (Aula~03). Na Figura~1 da
Aula~10, com $n=300$, o $1$-NN tem acurácia de treino $1{,}000$ e a pior acurácia de
teste das três, $0{,}826$, contra $0{,}888$ com $k=15$ e $0{,}889$ com $k=120$.
\end{solucao}

\end{document}
